\subsection{FORMATIVE CRITERIA}

CF1. If \(A\) and \(B\) are relations in a theory \(\mathfrak{G}\) , then \(\vee AB\) is a relation in \(\mathfrak{G}\) .

Consider two formative constructions (in \(\mathfrak{G}\) ), one of which contains \(A\) and the other \(B\) . Consider the sequence of assemblies obtained by writing first the assemblies of the first construction, then the assemblies of the second construction, and finally \(\vee AB\) . Since \(A\) and \(B\) are of the second species, it is immediately verified that this sequence is a formative construction of \(\mathfrak{G}\) . The assembly \(\vee AB\) is of the second species, hence it is a relation in \(\mathfrak{G}\) .

The three following criteria are established similarly:

CF2. If A is a relation in a theory C, then \(\neg A\) is a relation in C.

CF3. If A is a relation in a theory G, and if x is a letter, then \(\tau_{x}(A)\) is a term in G.

CF4. If \(A_1, A_2, \ldots, A_n\) are terms in a theory \(\mathfrak{C}\) , and if \(s\) is a relational (resp. substantific) sign of weight \(n\) in \(\mathfrak{C}\) , then \(sA_1A_2\ldots A_n\) is a relation (resp. a term) in \(\mathfrak{C}\) .

These criteria immediately imply the following :

CF5. If A and B are relations in a theory C, then \(\Longrightarrow AB\) is a relation in C.

CF6. Let \(A_1, A_2, \ldots, A_n\) be a formative construction in a theory \(\mathfrak{C}\) , and let \(x\) and \(y\) be letters. Suppose that \(y\) does not appear in any \(A_i\) . Then \((y|x)A_1\) , \((y|x)A_2, \ldots, (y|x)A_n\) is a formative construction in \(\mathfrak{C}\) .

To prove CF6, let \(A_i'\) be the assembly \((y|x)A_i\) . If \(A_i\) is a letter, then \(A_i'\) is a letter. If \(A_i\) is of the form \(\neg A_j\) , where \(A_j\) is an assembly of the second species which precedes \(A_i\) in the construction, then \(A_i'\) is identical with \(\neg A_j'\) by CS5, and \(A_j'\) is an assembly of the second species. The reasoning is similar if \(A_i\) is of the form \(\vee A_jA_k\) or \(sA_{j_1}A_{j_2}\ldots A_{j_m}\) , \(s\) being a specific sign of \(\mathfrak{G}\) . If finally \(A_i\) is of the form \(\tau_z(A_j)\) , where \(A_j\) is an assembly of the second species which precedes \(A_i\) ; in the construction, there are various cases to consider:

\begin{enumerate}
\def\labelenumi{(\alph{enumi})}
\item
  z is a letter distinct from x and y. Then \(A_{i}^{\prime}\) is identical with \(\tau_{z}(A_{j}^{\prime})\) by CS4, and \(A_{j}^{\prime}\) is an assembly of the second species.
\item
  z is identical with x. Then \(A_{i}\) does not contain x, hence \(A_{i}^{\prime}\) is identical with \(A_{i}\) , that is to say with \(\tau_{x}(A_{j})\) ; since y does not appear in \(A_{j}\) , \(\tau_{x}(A_{j})\) is identical with \(\tau_{y}(A_{j})\) by CS3.
\item
  z is identical with y. Then \(A_{i}\) is the assembly \(\tau A_{j}\) , because y does not appear in \(A_{j}\) ; therefore \(A_{i}^{\prime}\) is the assembly \(\tau A_{j}^{\prime}\) , that is \(\tau_{u}(A_{j}^{\prime})\) , where u is a letter which does not appear in \(A_{j}^{\prime}\) .
\end{enumerate}

CF7. Let \(A\) be a relation (resp. a term) in a theory \(\mathfrak{G}\) , and let \(x\) and \(y\) be letters. Then \((y|x)A\) is a relation (resp. a term) in \(\mathfrak{G}\) .

Let \(A_1, A_2, \ldots, A_n\) be a formative construction in which \(A\) appears. We shall show step by step that, if \(A_t\) is a relation (resp. a term) then \((y|x)A_t\) , which we shall denote by \(A_t'\) , is also a relation (resp. a term). Suppose that this point has been established for \(A_1, A_2, \ldots, A_{i-1}\) ; let us prove it for \(A_i\) . If \(A_i\) is a letter, then \(A_i'\) is a letter. If \(A_i\) is preceded in the construction by a relation \(A_j\) such that \(A_i\) is \(\neg A_j\) , then \(A_i'\) is identical with \(\neg A_j'\) by CS5, and \(\neg A_j'\) is a relation by CF2. The argument is similar if \(A_i\) is preceded by relations \(A_j, A_k\) such that \(A_i\) is \(\vee A_jA_k\) , or if \(A_i\) is preceded by terms \(A_{j_1}, \ldots, A_{j_m}\) such that \(A_i\) is \(sA_{j_1}\ldots A_{j_m}\) , where \(s\) is a specific sign of \(\mathfrak{G}\) of weight \(m\) . Finally, if \(A_i\) is preceded by a relation \(A_j\) such that \(A_i\) is \(\tau_z(A_j)\) , there are various cases to consider:

\begin{enumerate}
\def\labelenumi{(\alph{enumi})}
\item
  z is distinct from both x and y. Then \(A_{i}^{\prime}\) is identical with \(\tau_{\mathbf{z}}(A_{j}^{\prime})\) by CS4, and we know already that \(A_{j}^{\prime}\) is a relation; hence \(A_{i}^{\prime}\) is a term, by CF3.
\item
  \(z\) is identical with \(x\) . Then \(A_i\) does not contain \(x\) , therefore \(A_i'\) is identical with \(A_i\) , and consequently is a term.
\item
  \(z\) is identical with \(y\) . Then let \(u\) be a letter, distinct from both \(x\) and \(y\) , which does not appear in \(A_1, A_2, \ldots, A_j\) . By CF6, the sequence of assemblies \((u|y) A_1, \ldots, (u|y) A_j\) , which we shall denote by \(A_1'', \ldots, A_j''\) , constitutes a formative construction in \(\mathfrak{G}\) . Since \(y\) no longer appears in this new construction, \((y|x) A_1'', \ldots, (y|x) A_j''\) is a formative construction by CF6, so that \((y|x) A_j''\) is a relation in \(\mathfrak{G}\) ; consequently \(\tau_u((y|x)A_j'')\)
\end{enumerate}

is a term of \(\mathfrak{C}\) . But this term is identical with \((y|x)\tau_u(A_j'')\) by CS4, hence with \((y|x)\tau_y(A_j)\) by CS3, hence is identical with \(A_i'\) .

CF8. Let \(A\) be a relation (resp. a term) in a theory \(\mathfrak{G}\) , let \(x\) be a letter and \(T\) a term in \(\mathfrak{G}\) . Then \((T|x)A\) is a relation (resp. a term) in \(\mathfrak{G}\) .

Let \(A_1, A_2, \ldots, A_n\) be a formative construction in which \(A\) appears. Let \(x_1, x_2, \ldots, x_p\) be the distinct letters which appear in \(T\) . Let us associate with each letter \(x_i\) a letter \(x_i'\) , distinct from each of the letters \(x_1, \ldots, x_p\) and the letters which appear in \(A_1, \ldots, A_n\) , in such a way that the letters \(x_1', \ldots, x_p'\) are all distinct. The assembly

\[
\left(\boldsymbol {x} _ {1} ^ {\prime} \mid \boldsymbol {x} _ {1}\right) \left(\boldsymbol {x} _ {2} ^ {\prime} \mid \boldsymbol {x} _ {2}\right) \dots \left(\boldsymbol {x} _ {p} ^ {\prime} \mid \boldsymbol {x} _ {p}\right) \boldsymbol {T}
\]

is a term \(\pmb{T}'\) by CF7, and \((\pmb{T}|\pmb{x})\pmb{A}\) is identical with

\[
\left(\boldsymbol {x} _ {1} \mid \boldsymbol {x} _ {1} ^ {\prime}\right) \left(\boldsymbol {x} _ {2} \mid \boldsymbol {x} _ {2} ^ {\prime}\right) \dots \left(\boldsymbol {x} _ {p} \mid \boldsymbol {x} _ {p} ^ {\prime}\right) \left(\boldsymbol {T} ^ {\prime} \mid \boldsymbol {x}\right) \boldsymbol {A}
\]

by application of CS1. It is therefore enough to show that \((T^{\prime}|x)A\) is a relation (resp. a term); in other words, we may suppose from now on that the letters which appear in \(T\) do not appear in \(A_{1},\ldots ,A_{n}\) . We shall show step by step that, if \(A_{t}\) is a relation (resp. a term), then \((T|x)A_{t}\) , which we shall denote by \(A_{t}^{\prime}\) , is a relation (resp. a term). Suppose this point has been established for \(A_{1},A_{2},\ldots ,A_{i - 1}\) , and let us prove it for \(A_{i}\) . If \(A_{i}\) is a letter, then \(A_{i}^{\prime}\) is either the same letter or \(T\) , and therefore a term. If \(A_{i}\) is of the form \(\lnot A_j\) , where \(A_{j}\) is a relation which precedes \(A_{i}\) in construction, then \(A_{i}^{\prime}\) is identical with \(\lnot A_j^{\prime}\) by CS5, and we know already that \(A_{j}^{\prime}\) is a relation, hence \(A_{i}^{\prime}\) is a relation by CF2. The proof is analogous if \(A_{i}\) is of the form \(\vee A_{j}A_{k}\) , or \(sA_{j_1}\ldots A_{j_m}\) . Finally, if \(A_{i}\) is of the form \(\tau_{z}(A_j)\) , where \(A_{j}\) is a relation which precedes \(A_{i}\) in the construction, there are various cases to be considered:

\begin{enumerate}
\def\labelenumi{(\alph{enumi})}
\item
  \(z\) is distinct from \(x\) and from the letters which appear in \(T\) . Then \(A_i'\) is identical with \(\tau_z(A_j')\) by CS4, and we know already that \(A_j'\) is a relation; hence \(A_i'\) is a term by CF3.
\item
  \(\pmb{z}\) is identical with \(\pmb{x}\) . Then \(A_{i}\) does not contain \(\pmb{x}\) , hence \(A_{i}'\) is identical with \(A_{i}\) and consequently is a term.
\item
  \(z\) appears in \(T\) . Then \(z\) does not appear in \(A_j\) , so that \(A_i'\) is identical with \(\tau A_j\) ; hence \(A_i'\) is identical with \(\tau A_j'\) . Now, we know already that \(A_j'\) is a relation, and \(\tau A_j'\) is identical with \(\tau_u(A_j')\) , where \(u\) is a letter which does not appear in \(A_j'\) ; it follows by CF3 that \(A_i'\) is a term.
\end{enumerate}

Intuitively, if A is a relation in C, which we may regard as expressing a property of an object x, the assertion \((B|x)A\) amounts to saying that the object B has this property. If A is a term in C, it represents an object which depends in some way on the object denoted by x; the term \((B|x)A\) represents what the object A becomes when we take x to be the object B.

